% Chronophotographie d'un palet sur une table à coussin d'air, échelle 1 (exercice 1, chapitre 13)
\begin{tikzpicture}[x=1cm,y=1cm,rotate=8,
  vit/.style={otNavy,line width=1pt,-{Stealth[length=2mm]}}]
\foreach \i in {0,...,6}{%
  \draw[otInk!70,line width=0.5pt] ({1.6*\i},0) circle (0.13);
  \fill[otInk] ({1.6*\i},0) circle (0.035);
  \node[pcLegende,anchor=north] at ({1.6*\i},-0.18) {$M_{\i}$};}
\draw[otInk!70,line width=0.4pt,{Stealth[length=1.3mm]}-{Stealth[length=1.3mm]}] (0,0.55) -- (1.6,0.55) node[midway,pcLegende,anchor=south] {\qty{1,6}{cm}};
\solution{
  \draw[vit] (4.8,0.3) -- ++(4.0,0) node[pcLegende,anchor=south,text=otNavy] {$\vec{v}_3$};
  \draw[vit] (6.4,-0.75) -- ++(4.0,0) node[pcLegende,anchor=north,text=otNavy] {$\vec{v}_4$};}
\end{tikzpicture}
\par\smallskip\legende{Positions du centre du palet toutes les \qty{40}{ms} (document à l'échelle 1).\solution{ Vecteurs vitesse à l'échelle \qty{1}{cm} pour \qty{0,10}{m.s^{-1}} : $\vec{v}_4=\vec{v}_3$.}}
